\apisummary{
  Collectively allocate symmetric memory in a specific space.
}

\begin{apidefinition}

\begin{Csynopsis}
void *@\FuncDecl{shmem\_space\_malloc}@(shmem_space_t space, size_t size);
\end{Csynopsis}

\begin{apiarguments}
  \apiargument{IN}{space}{The space within which to allocate.}
  \apiargument{IN}{size}{The size, in bytes, of a block to be
    allocated from the space.}
\end{apiarguments}


\apidescription{
  The \FUNC{shmem\_space\_malloc} routine is a collective operation on the
  team associated with \VAR{space} and returns the symmetric address of a
  block of at least \VAR{size} bytes, which shall be suitably aligned
  so that it may be assigned to a pointer to any type of object.
  This block is allocated from the specific symmetric heap identified
  by \VAR{space}.
  When \VAR{size} is zero, the \FUNC{shmem\_space\_malloc} routine performs
  no action and returns a null pointer; otherwise,
  \FUNC{shmem\_space\_malloc} calls a procedure that is semantically equivalent 
  to \FUNC{shmem\_team\_sync} on exit. This ensures
  that all \acp{PE} with access to the space participate
  in the memory allocation, and that the object on other \acp{PE} can be used as soon as the local
  \ac{PE} returns. 
  The values of the \VAR{space} and \VAR{size} arguments must be identical on all
  \acp{PE}; otherwise, the behavior is undefined.
}

\apireturnvalues{
  The \FUNC{shmem\_space\_malloc} routine returns the symmetric address of
  the allocated space; otherwise, it returns a null pointer.
}

\end{apidefinition}
